\begin{pastexam}[label={px:ccu-109-swa-08}]{中正資工所甲組 109 學年度 第 13 題}
（是非改錯）

(18\%) True or false. If the statement is false, correct the wrong part. Simply negate the statement
is not accepted. (3\% each)
\begin{enumerate}[label=(\arabic*)]
\item Given an array A[1\,..\,n] of integers, the running time of Counting Sort is polynomial in the
input size n.
\item Given an array A[1\,..\,n] of integers, the running time of Heap Sort is polynomial in the input
size n.
\item Depth-first search will take $\Theta$(V$^2$) time on a graph G = (V, E) represented as an
adjacency matrix.
\item For every dynamic program, we can assign weights to edges in the directed acyclic graph of
dependences among subproblems, such that finding a shortest path in this DAG is equivalent to
solving the dynamic program.
\item If a problem X can be reduced to a known NP-hard problem, then X must be NP-hard.
\item Using DFS to search augmenting paths in Ford-Fulkerson algorithm could reduce the time
complexity
\end{enumerate}
\end{pastexam}

\begin{pastsoln}
答案：(1) F；(2) T；(3) T；(4) F；(5) F；(6) F。

\textbf{(1) F}。 計數排序的時間是 $\Theta(n + k)$，$k$ 是數值的範圍（資結書第~15.2.1~節）。 題目
只說是整數，沒有限制大小，$k$ 可以遠大於 $n$。 例如 $n = 2$，$A = \langle 0, 2^m \rangle$：輸入只要
大約 $m$ 個位元就寫得下，計數排序卻要開 $2^m + 1$ 格的計數陣列，時間是輸入長度的指數。 這是
擬多項式時間（資結書第~1.9、18.5.5~節），不是多項式。 \textbf{改正}：the running time of Counting
Sort is $\Theta(n + k)$, polynomial in $n$ and the value range $k$ (pseudo-polynomial), \textbf{not}
polynomial in the input size。

\textbf{(2) T}。 堆積排序做 $O(n \lg n)$ 次比較和交換（資結書第~14.6~節），次數和數值大小無關。
每個整數至少佔一個位元，輸入長度 $N \ge n$，所以 $n \lg n \le N^2$。 就算把比較兩個 $b$ 位元整數
的成本算成 $O(b)$，總時間 $O(n \lg n \cdot b)$ 也不超過 $O(N^2)$，仍是輸入長度的多項式。

\textbf{(3) T}。 用鄰接矩陣時，DFS 每個頂點恰好拜訪一次（資結書虛擬碼~21.3 的外層迴圈會從每個
還沒拜訪的頂點出發），拜訪頂點 $u$ 時要把矩陣第 $u$ 列的 $V$ 格全部看過才知道誰是鄰居。 每列
看一次，共 $V \cdot V$ 格，再加上外層迴圈的 $\Theta(V)$，時間 $\Theta(V^2)$，和邊數無關
（資結書第~21.2、21.4~節）。 鄰接串列才是 $\Theta(V + E)$。

\textbf{(4) F}。 化成 DAG 最短路徑，要求每個子問題的遞迴式是「在幾個選擇裡取最小，每個選擇只用到
\textbf{一個}子問題，再加上一個固定的成本」，這樣成本才能放在邊上，一條路徑只經過一個前驅。
很多動態規劃不是這個形狀：
\begin{itemize}
\item 矩陣鏈乘積 $m[i, j] = \min_{i \le k < j} \bigl(m[i, k] + m[k+1, j] + p_{i-1} p_k p_j\bigr)$
（資結書第~18.2~節）：每個選擇同時用到 $m[i, k]$ 和 $m[k+1, j]$ 兩個子問題，相依關係是「一個點
接兩個點」的超圖，一條路徑表示不了。 最佳二元搜尋樹（資結書第~18.4~節）也一樣。
\item 計數型的動態規劃，例如 $n$ 個節點的二元樹個數 $b_n = \sum_k b_k b_{n-1-k}$，要的是加總和乘積，
根本不是取最小。
\end{itemize}
\textbf{改正}：For \textbf{some} dynamic programs (those whose recurrence takes a minimum over choices,
each depending on a single subproblem plus a cost), we can assign weights \ldots。 這一類（例如 0/1
背包、LCS、編輯距離）才能看成 DAG 上的最短或最長路徑。

\textbf{(5) F}。 歸約的方向反了。 「$X$ 可以歸約到 $Y$」（$X \le_p Y$）表示 $Y$ 至少和 $X$ 一樣難，
只能推出 $X$ 不比 $Y$ 難（資結書第~28.3~節）。 反例：$X$ 是「不管輸入是什麼，答案都是是」的問題。
把每個輸入都轉成同一個可滿足的公式（例如 $x_1$），就是 $X$ 到 SAT 的多項式歸約，SAT 是 NP 難的。
但 $X$ 不是 NP 難：SAT 有答案為否的輸入，要歸約到 $X$ 就得把它轉成 $X$ 答案為否的輸入，而 $X$ 沒有
這種輸入，所以 SAT 不能歸約到 $X$。 \textbf{改正}：If a known NP-hard problem can be reduced to
problem X (in polynomial time), then X must be NP-hard。 理由是歸約可以串接（資結書引理~28.3）：
NP 裡任何問題 $\le_p$ 已知的 NP 難問題 $\le_p X$。

\textbf{(6) F}。 Ford--Fulkerson 不規定怎麼找增廣路徑，整數容量時的界是 $O(E\,|f^*|)$，$|f^*|$ 是
最大流量（資結書第~26.3~節）。 用 DFS 找路徑並沒有改進這個界：資結書第~26.3.3~節的例子裡，DFS
可能每次都走經過中間那條容量 $1$ 的邊，一次只送 $1$，輪數和容量成正比。 能降低時間的是用
\textbf{BFS} 找邊數最少的增廣路徑，也就是 Edmonds--Karp 演算法，時間 $O(VE^2)$，和容量無關
（資結書第~26.4~節、定理~26.3）。 \textbf{改正}：Using \textbf{BFS} to search augmenting paths
(the Edmonds--Karp algorithm) could reduce the time complexity to $O(VE^2)$。
\end{pastsoln}
